\documentclass[11pt]{article}
\usepackage[margin=1in]{geometry}
\usepackage[T1]{fontenc}
\usepackage[utf8]{inputenc}
\usepackage{amsmath,amssymb,mathtools}
\usepackage{booktabs}
\usepackage{hyperref}
\usepackage{xcolor}
\usepackage{listings}

\hypersetup{
  colorlinks=true,
  linkcolor=blue!60!black,
  urlcolor=blue!60!black
}

\definecolor{codebg}{RGB}{248,248,248}
\definecolor{codeborder}{RGB}{210,210,210}
\definecolor{codecomment}{RGB}{80,120,80}
\definecolor{codekeyword}{RGB}{30,60,150}

\lstdefinestyle{reportcode}{
  backgroundcolor=\color{codebg},
  basicstyle=\ttfamily\small,
  breaklines=true,
  columns=fullflexible,
  commentstyle=\color{codecomment},
  frame=single,
  framerule=0.5pt,
  keywordstyle=\color{codekeyword},
  rulecolor=\color{codeborder},
  keepspaces=true,
  showstringspaces=false,
  tabsize=2
}

\lstset{style=reportcode, language=Python}

\setlength{\parindent}{0pt}
\setlength{\parskip}{0.5em}

\newcommand{\R}{\mathbb{R}}
\newcommand{\ysp}{y_{\mathrm{sp}}}
\newcommand{\uvec}{u}
\newcommand{\xvec}{x}
\newcommand{\Aaug}{A_{\mathrm{aug}}}
\newcommand{\Baug}{B_{\mathrm{aug}}}
\newcommand{\Caug}{C_{\mathrm{aug}}}
\newcommand{\Qout}{Q_{\mathrm{out}}}
\newcommand{\Rin}{R_{\mathrm{in}}}
\newcommand{\Hp}{H_p}
\newcommand{\Hc}{H_c}
\newcommand{\argmin}{\mathop{\mathrm{arg\,min}}}


\title{Residual Methods Under Model Mismatch}
\date{2026-04-01}

\begin{document}
\maketitle

\section*{Method Goal}

The residual notebook family adds a second control channel on top of baseline
MPC. The MPC still computes a feasible baseline move, but a TD3 policy proposes
an extra residual move that is clipped by mismatch-aware safety bounds before it
is applied. The intent is to preserve the optimization structure of MPC while
letting RL compensate for mismatch that the linear model does not capture well.

\section*{Notebook Sources Covered}

\begin{itemize}
\item \texttt{RL\_assisted\_MPC\_residual\_model\_mismatch.ipynb}
\item \texttt{RL\_assisted\_MPC\_residual\_model\_mismatch1.ipynb}
\item \texttt{RL\_assisted\_MPC\_residual\_model\_mismatch2.ipynb}
\item \texttt{RL\_assisted\_MPC\_residual\_model\_mismatch\_multi.ipynb}
\end{itemize}

\section*{State, Action, Reward, And Update Logic}

The residual family augments the standard scaled state with innovation and
tracking error terms:
\[
s_t =
\begin{bmatrix}
\tilde{x}_{t} \\
\tilde{y}_{\mathrm{sp},t} \\
\tilde{u}_{t}^{\mathrm{dev}} \\
\mathrm{innov}_t \\
e^{\mathrm{track}}_t
\end{bmatrix},
\]
where
\[
\mathrm{innov}_t = y_t^{\mathrm{scaled}} - \Caug \hat{x}_t,
\qquad
e^{\mathrm{track}}_t = y_t^{\mathrm{scaled}} - \ysp^{(t)}.
\]

The agent proposes a residual move $\Delta u_t^{\mathrm{res}}$ in scaled space.
That proposal is not applied directly. It is clipped by two intersected bounds:

\begin{itemize}
\item a magnitude bound relative to the MPC move,
\item a safety bound given by actuator headroom around the baseline input.
\end{itemize}

The applied control is therefore
\[
u_t^{\mathrm{applied}} = u_t^{\mathrm{base}} + \Delta u_t^{\mathrm{res,exec}},
\]
where $\Delta u_t^{\mathrm{res,exec}}$ is the clipped residual that actually
survives the bound intersection.

\section*{Mismatch Bands And Residual Gating}

The notebook family defines mismatch bands in scaled output coordinates:
\[
\mathrm{band}_{\mathrm{scaled}} =
\frac{\max(k_{\mathrm{rel}} |\ysp^{\mathrm{phys}}|,
\mathrm{band\_floor}^{\mathrm{phys}})}
{\mathrm{data\_max} - \mathrm{data\_min}}.
\]
It also defines soft and hard gating helpers that estimate how far the current
tracking error is outside the allowed band.

\begin{lstlisting}
def compute_band_scaled(y_sp_phys, data_min, data_max, n_inputs,
                        k_rel, band_floor_phys):
    dy = np.maximum(data_max[n_inputs:] - data_min[n_inputs:], 1e-12)
    band_phys = np.maximum(k_rel * np.abs(y_sp_phys), band_floor_phys)
    return band_phys / dy

def residual_outside_weight(e_track_scaled, band_scaled,
                            tau_frac=0.7, kind="geom"):
    abs_e = np.abs(e_track_scaled)
    tau = np.maximum(tau_frac * band_scaled, 1e-12)
    x = (band_scaled - abs_e) / tau
    s = 1.0 / (1.0 + np.exp(-x))
    ...
    return 1.0 - w_in
\end{lstlisting}

\section*{Core Residual Control Loop}

\begin{enumerate}
\item Compute the baseline MPC move from the current observer state.
\item Compute innovation, tracking error, and a mismatch-band measure.
\item Ask the TD3 policy for a raw residual move proposal.
\item Build a magnitude envelope around the baseline MPC move.
\item Intersect that envelope with actuator headroom bounds.
\item Clip the residual move, apply the corrected control, and log the executed
      residual rather than the raw proposal.
\item Store the executed residual action in replay so the TD3 buffer reflects
      what the plant actually saw.
\end{enumerate}

\section*{Key Code Paths}

\subsection*{Notebook-local residual safety logic}

\begin{lstlisting}
u_base = sol.x[:n_inputs] + ss_scaled_inputs
u_base = np.clip(u_base, u_min_scaled_abs, u_max_scaled_abs)

delta_u_mpc = (u_base - u_current).astype(np.float32)
delta_u_res_raw = map_to_bounds(a_raw, low_res, high_res).astype(np.float32)

mag = (rho * beta_res) * (np.abs(delta_u_mpc) + du0_res)
low_mag = -mag
high_mag = mag

low_safe = (u_min_scaled_abs - u_base).astype(np.float32)
high_safe = (u_max_scaled_abs - u_base).astype(np.float32)

low_bound = np.maximum(low_mag, low_safe)
high_bound = np.minimum(high_mag, high_safe)
delta_u_res = np.clip(delta_u_res_raw, low_bound, high_bound)
\end{lstlisting}

\subsection*{Notebook-local replay action logic}

\begin{lstlisting}
u_applied_scaled_abs = u_base + delta_u_res
u_applied_scaled_abs = np.clip(u_applied_scaled_abs,
                               u_min_scaled_abs, u_max_scaled_abs)

delta_u_res_exec = (u_applied_scaled_abs - u_base).astype(np.float32)
a_exec = map_from_bounds(delta_u_res_exec, low_res, high_res).astype(np.float32)
a_exec = np.clip(a_exec, -1.0, 1.0).astype(np.float32)
\end{lstlisting}

\section*{Variants Covered}

\begin{itemize}
\item The base residual mismatch notebook defines the main residual scaffold.
\item Variants \texttt{1} and \texttt{2} keep the same pattern while changing
      disturbance or mismatch schedules.
\item The \texttt{multi} variant extends the idea to a more involved residual
      setting while preserving the same core ingredients: innovation terms,
      mismatch bands, clipped residual execution, and replay of executed moves.
\end{itemize}

\section*{Limitations, Caveats, And Open Questions}

\begin{itemize}
\item This family depends on several notebook globals such as
      \texttt{TEST\_CYCLE}, \texttt{IC\_opt}, \texttt{bnds}, \texttt{cons},
      \texttt{k\_rel}, and \texttt{band\_floor\_phys}.
\item Residual replay stores the executed clipped action, not the raw action.
      That is intentional and must be preserved in refactors.
\item Additional mismatch scheduling includes changes to \texttt{CMf}, not just
      the shared $Q_i$, $Q_s$, and $hA$ schedules.
\item This method is explicitly notebook-local and should not be inferred from
      stale helper modules.
\end{itemize}

\end{document}
